\documentclass[12pt]{article}

\usepackage[a4paper, total={6in, 8in}]{geometry}
\usepackage{textcomp}

\begin{document}
\setlength{\parindent}{0pt}

\def \t {\textrightarrow}
\def \v {\vspace{1em}}
\def \bi {\begin{itemize}}
\def \ei {\end{itemize}}
\def \s[#1] {\section*{#1}}
\def \ss[#1] {\subsection*{#1}}
\def \sss[#1] {\subsubsection*{#1}}

\s[D'Annunzio]
Scappa a napoli, dove trova un ambiente che lo accoglie \t sempre in un ambiente di vita libera, con duelli e gioco \\
Pero lui ama l'arte, come possibilita di vita \t arte e vita ("Letteratura come vita" è il titolo di un documento di Carlo Bò) \\
In D'annuznio abbiamo totale sovrapposizione tra arte e vita \t non c'è spaccatura tra d'annunzio letterato e d'annunzio uomo \\
Questo permette di agganciarsi a "Il piacere", con il personaggio "Andreas Perelli" del 1889 \\
1889: data rilevante, perche si erano gia verificati dei fatti importanti (letterari) \t sono anche gli anni letterari di Verga, e visione sinottica con istanze differenti che danno opere differenti, ma sono tutti contemporanei nella fase della crisi dell'io alla fine del secolo \\
Perelli è d'annunzio trapiantato in un opera scritta, circondata di oggetti e belezza \t trae vita da un binomio che è la letteratura e la vita, ma proprio perchè nella morale decadente si vedono i parametri del bene e del male, qua la morale non è una morale esterna ma insita al personaggio \\
Si assiste a una autonomia dell'arte dalla morale tradizionale, che viene espresso dalla filosofia di Benedetto Croce (di marcatura idealista e fautore di una certa visione della storia della critica, la critica crociana, che è diversa da quella conosciuta con foscolo) \\
La critica foscoliana è di marcatura + storica, mentre quella di crosce + filologica e contestuale \\
Aleggia sempre il decadentismo, di cui D'annunzio è comunque vassallo

\ss[Il piacere]
Qua c'è una morale non morale \\
Qua la morale è nella non-morale dell'opera d'arte \t che è di per se morale \\
Tutto ciò che è al di fuori è sovrastruttura \\
Le donne hanno molte sfaccettature: è rassicurante e materna, che è Maria che è importante per perelli \t ma poi c'è Elena che è la seduzione \t si sviluppa in una malattia morbosa che si sviluppa nel personaggio, che è continuamente coinvolto eroticamente \\
Vive la sessualita in modo istintivo e brado \\
L'ero estetizzante è pero solo \t come anche Dorian Gray \t l'eroe è solo, ed è disprezzante della massa \\
Tutta la sua fortuna mediatica è concentrata sull'imagine di perelli, che è il suo alter ego \t con solitudine, duelli, alsdkfj \\
Alla fine Elena lo lascia \t eros e tanatos si chiudono nel loro ciclo \t per la sofferenza del vuoto, dell'horror vacui \\
Si chiude con l'immagine di andreas perelli che si trova solo e ha perso tutto, perchè in fondo non ha mai avuto se stesso \\
Questa è la parabola di un eroe decaente \t da collegare con Hyusman con "A reboir", e "The picture of Dorian Gray" \\
Per quanto riguarda il linguaggio, colleghiamo anche il linguaggio elettivo del Parnassianesimo (dal monte parmaso, dove si trovano le muse)

\v

D'annunzio è anche autore di tragedie macabre, come (vedere nome, citta morta) \\
Ha un animo insaziabile \\
Non è presente alcuna idealizzazione e alcuna donna angelo \t perche la idealizzazione prevede anche una ?? \t e c'è attenenza al superuomo nietzschiano

\v

L'autore ha una passione per la musica di Wagner \t perche ha sentimenti forti \t il personaggio Perelli va anche al funerale di Wagner ??? \\
"Le vergini delle rocce" \t vicinanza a Nietzsche: matrimonio tra Claudio Cantelmo e una ninfa \t che sceglie per dare vita al futuro superuomo \\
Cantelmo è uno scultore \t e un figlio non lo puo scolpire \\
L'autore aveva figli d'appertutto \t sara monogamo per un periodo solo quando incontrera Aduse ??
\end{document}
